\documentclass[10pt,aspectratio=169]{beamer}
% Preámbulo modular para presentaciones Beamer (Modelación Estadística)
% Diseño y paleta institucional del Tecnológico de Monterrey

\documentclass[aspectratio=169,xcolor={dvipsnames,table}]{beamer}

\usepackage[utf8]{inputenc}
\usepackage[spanish,mexico]{babel}
\usepackage{amsmath,amssymb,amsfonts,mathtools}
\usepackage{graphicx}
\usepackage{listings}
\usepackage{booktabs}
\usepackage{tikz}
\usetikzlibrary{matrix,arrows,decorations.pathmorphing,shapes.geometric,positioning}

% Paleta Institucional Tecnológico de Monterrey
\definecolor{TecAzul}{HTML}{0039A6}       % Azul TEC: color primario institucional
\definecolor{TecAzulOscuro}{HTML}{1A2E51} % Azul marino oscuro
\definecolor{TecRojo}{HTML}{EC2661}       % Rojo de acento (paleta miTec)
\definecolor{TecGris}{HTML}{646464}       % Gris neutro para comentarios y subtítulos
\definecolor{TecFondoBloque}{HTML}{F4F6F9}% Fondo suave para bloques

% Configuración de Tema Beamer
\usetheme{Boadilla}
\usecolortheme[named=TecAzulOscuro]{structure}
\setbeamercolor{alerted text}{fg=TecRojo}
\setbeamercolor{palette primary}{bg=TecAzulOscuro,fg=white}
\setbeamercolor{palette secondary}{bg=TecAzul,fg=white}
\setbeamercolor{palette tertiary}{bg=TecAzulOscuro!80!black,fg=white}
\setbeamercolor{title}{fg=TecAzulOscuro,bg=TecFondoBloque}
\setbeamercolor{frametitle}{fg=TecAzulOscuro,bg=TecFondoBloque}

% Bloques personalizados elegantes
\setbeamercolor{block title}{bg=TecAzulOscuro,fg=white}
\setbeamercolor{block body}{bg=TecFondoBloque,fg=black}
\setbeamercolor{block title alerted}{bg=TecRojo,fg=white}
\setbeamercolor{block body alerted}{bg=TecRojo!8,fg=black}
\setbeamercolor{block title example}{bg=TecAzul,fg=white}
\setbeamercolor{block body example}{bg=TecAzul!8,fg=black}

% Redondear bordes y añadir sombra sutil
\setbeamertemplate{blocks}[rounded][shadow=true]
\setbeamertemplate{navigation symbols}{}

% Pie de página limpio con número de diapositiva
\setbeamertemplate{footline}{
  \leavevmode%
  \hbox{%
  \begin{beamercolorbox}[wd=.7\paperwidth,ht=2.25ex,dp=1ex,left]{author in head/foot}%
    \hspace*{2ex}\insertshortauthor~~(\insertshortinstitute) \hfill \insertshorttitle
  \end{beamercolorbox}%
  \begin{beamercolorbox}[wd=.3\paperwidth,ht=2.25ex,dp=1ex,right]{date in head/foot}%
    \insertshortdate{}\hspace*{2em}
    \insertframenumber{} / \inserttotalframenumber\hspace*{2ex}
  \end{beamercolorbox}}%
  \vskip0pt%
}

% Configuración de Listings de Python para visualización pedagógica de código
\lstset{
    language=Python,
    basicstyle=\ttfamily\footnotesize,
    keywordstyle=\color{TecAzulOscuro}\bfseries,
    stringstyle=\color{TecRojo},
    commentstyle=\color{TecGris}\itshape,
    showstringspaces=false,
    breaklines=true,
    frame=lines,
    rulecolor=\color{TecAzulOscuro!30},
    backgroundcolor=\color{TecFondoBloque},
    aboveskip=0.5em,
    belowskip=0.5em,
    literate={á}{{\'a}}1 {é}{{\'e}}1 {í}{{\'i}}1 {ó}{{\'o}}1 {ú}{{\'u}}1
             {Á}{{\'A}}1 {É}{{\'E}}1 {Í}{{\'I}}1 {Ó}{{\'O}}1 {Ú}{{\'U}}1
             {ñ}{{\~n}}1 {Ñ}{{\~N}}1 {¿}{{?`}}1 {¡}{{!`}}1
}

% Metadatos institucionales por defecto
\title[Modelación Estadística]{Modelación Estadística para Ciencia de Datos e Ingeniería}
\author[J. Castillo Colmenares]{Juliho David Castillo Colmenares}
\institute[Tec de Monterrey]{Tecnológico de Monterrey \\ Escuela de Ingeniería y Ciencias}
\date{\today}
% Comandos y macros estadísticas compartidas para las presentaciones Beamer (Metropolis)

% Conjuntos numéricos básicos
\newcommand{\R}{\mathbb{R}}
\newcommand{\N}{\mathbb{N}}
\newcommand{\Q}{\mathbb{Q}}
\newcommand{\Z}{\mathbb{Z}}
\newcommand{\set}[1]{\left\{ #1 \right\}}
\newcommand{\abs}[1]{\left|#1\right|}
\newcommand{\norm}[1]{\left\| #1 \right\|}

% Comandos estadísticos y probabilísticos del curso
\newcommand{\Var}{\operatorname{Var}}
\newcommand{\cov}{\operatorname{Cov}}
\newcommand{\corr}{\rho}
\newcommand{\comb}[2]{\begin{pmatrix} #1 \\ #2 \end{pmatrix}}
\newcommand{\s}{\sigma}
\newcommand{\particion}{\mathfrak{P}}
\newcommand{\card}[1]{\#\left( #1 \right)}
\newcommand{\seq}[3]{\set{#1}_{#2}^{#3}}
\newcommand{\E}{\mathbb{E}}
\newcommand{\Prob}{\mathbb{P}}

% Entornos pedagógicos y cajas en español académico natural
\newenvironment{discusion}{%
  \begin{alertblock}{Pregunta de Discusión}%
}{%
  \end{alertblock}%
}

\newenvironment{reflexion}{%
  \begin{alertblock}{Reflexión para el Aula}%
}{%
  \end{alertblock}%
}

\newenvironment{paradoja}{%
  \begin{alertblock}{Paradoja de Exploración}%
}{%
  \end{alertblock}%
}

\newenvironment{motivacion}{%
  \begin{exampleblock}{Motivación: Ciencia de Datos en la Práctica}%
}{%
  \end{exampleblock}%
}

\newenvironment{retoclase}{%
  \begin{block}{Reto Rápido en Equipo}%
}{%
  \end{block}%
}

\title[Introducción a la regresión lineal]{Sección 09.02: Introducción a la regresión lineal}\subtitle{Modelos, supuestos y aplicaciones}\author[J. Castillo Colmenares]{Juliho Castillo Colmenares}\institute{Tecnológico de Monterrey}\date{}
\begin{document}
\begin{frame}[plain]\titlepage\end{frame}
\begin{frame}{Hoja de ruta}\small\begin{enumerate}\item Contexto y modelos.\item Correlación, supuestos y tipos.\item Aplicaciones y flujo de trabajo.\end{enumerate}\end{frame}
\begin{frame}{Fuente y alcance}\small La nota introduce la regresión lineal como técnica predictiva y prepara las secciones matemáticas y computacionales del capítulo.\begin{block}{Pregunta guía}¿Qué hace que una relación lineal sea un modelo estadístico útil?\end{block}\end{frame}
\begin{frame}{Regresión lineal}\small La regresión construye una relación funcional entre variables explicativas y una respuesta, usando datos históricos para estimar parámetros y producir predicciones. Es un modelo estocástico: incorpora un error aleatorio que distingue las observaciones reales de una relación determinística.\end{frame}
\begin{frame}{Usos principales}\small\begin{itemize}\item Predicción de valores futuros.\item Inferencia sobre asociaciones.\item Control de procesos.\item Investigación científica.\end{itemize}\end{frame}
\begin{frame}{Ejemplo de estructura}\small Para el precio $P$ de una casa con tamaño $S$, comodidades $A$ y transporte $T$,\[P=a_1S+a_2A+a_3T+\epsilon.\]\end{frame}
\begin{frame}{Parámetros y variables}\small $P$ es la respuesta; $S,A,T$ son entradas conocidas; $a_1,a_2,a_3$ son parámetros que deben estimarse; $\epsilon$ recoge factores no observados.\end{frame}
\begin{frame}{Estimación}\small Los parámetros se obtienen a partir de datos históricos. Una vez estimados, el modelo puede producir valores ajustados y predicciones nuevas.\end{frame}
\begin{frame}{Correlación y regresión}\small\begin{itemize}\item Correlación: asociación lineal simétrica y sin unidades.\item Regresión: relación direccional para modelar y predecir.\end{itemize}\end{frame}
\begin{frame}{Cuidado causal}\small Una correlación alta puede sugerir un ajuste lineal, pero no demuestra causalidad ni sustituye un diseño experimental.\end{frame}
\begin{frame}{Supuesto de linealidad}\small La media condicional de la respuesta debe cambiar linealmente con los predictores, al menos en el rango de interés.\end{frame}
\begin{frame}{Supuesto de independencia}\small Las observaciones deben ser independientes. Dependencia temporal, espacial o por agrupamiento requiere un modelo o un diseño distinto.\end{frame}
\begin{frame}{Supuesto de homocedasticidad}\small La varianza de los errores debe ser aproximadamente constante a través de los niveles de los predictores.\end{frame}
\begin{frame}{Supuesto de normalidad}\small La normalidad de errores es especialmente relevante para inferencia en muestras pequeñas; no es una condición de toda predicción.\end{frame}
\begin{frame}{No multicolinealidad}\small En regresión múltiple, predictores muy correlacionados dificultan separar sus efectos y vuelven inestables los coeficientes.\end{frame}
\begin{frame}{Regresión simple}\small Con un predictor,\[Y=\beta_0+\beta_1X+\epsilon.\]La pendiente resume el cambio esperado de $Y$ por unidad de $X$.\end{frame}
\begin{frame}{Regresión múltiple}\small Con varios predictores,\[Y=\beta_0+\beta_1X_1+\cdots+\beta_pX_p+\epsilon.\]Cada coeficiente es un efecto parcial condicionado a los demás.\end{frame}
\begin{frame}{Mínimos cuadrados}\small Ambos tipos estiman parámetros minimizando la suma de errores cuadráticos. La elección de variables debe respetar la pregunta y los supuestos.\end{frame}
\begin{frame}{Aplicaciones}\small Finanzas, marketing, medicina, ingeniería y ciencias sociales usan regresión para cuantificar asociaciones y generar predicciones, siempre con validación.\end{frame}
\begin{frame}{Flujo de trabajo}\small\begin{enumerate}\item Definir respuesta y predictores.\item Ajustar.\item Diagnosticar supuestos.\item Validar fuera de muestra.\item Comunicar incertidumbre.\end{enumerate}\end{frame}
\begin{frame}{Síntesis}\small La regresión lineal es un modelo estocástico: estima una relación, separa señal de ruido y exige supuestos explícitos para interpretar sus resultados.\end{frame}
\begin{frame}{Conexión con el capítulo}\small Las siguientes secciones desarrollan las matemáticas, implementaciones, validación y extensiones de regresión lineal.\end{frame}
\end{document}
